\section{Introduction}
\label{sec:introduction}

German Whist is a two-player trick-taking card game that presents a fascinating
blend of imperfect and perfect information. Unlike many card games where
trick-taking is the sole objective, German Whist divides play into two distinct
phases: a card-acquisition phase where players compete for face-up cards to
build their hands, followed by a standard trick-taking phase where the majority
of tricks won determines the winner.

This two-phase structure creates a rich strategic landscape. In Phase~1, players
must balance immediate tactical considerations (winning or losing each trick)
against long-term strategic goals (assembling a strong hand for Phase~2). The
information asymmetry---the winner sees and takes the face-up stock card while
the loser receives an unknown card---adds a further dimension of complexity.

\subsection{Motivation}

Our goal is twofold:
\begin{enumerate}
  \item \textbf{Build a strong AI opponent} capable of playing German Whist at
    a high level, combining game-tree search for the endgame with probabilistic
    reasoning for the opening.
  \item \textbf{Identify optimal strategic principles} through large-scale
    simulation, quantifying which factors most influence game outcomes.
\end{enumerate}

\subsection{Approach}

We decompose the problem by game phase:

\begin{itemize}
  \item \textbf{Phase~2 (perfect information):} Both players' hands are
    deducible through card counting. We solve this exactly using alpha-beta
    minimax with transposition tables, achieving optimal play in under one
    second for all 13-card endgames.
  \item \textbf{Phase~1 (imperfect information):} The opponent's hand and stock
    order are unknown. We use a determinized evaluation approach: sampling
    plausible opponent hands, simulating trick outcomes, and evaluating the
    resulting hand quality.
\end{itemize}

We validate our AI through a tournament framework, pitting it against random
and heuristic opponents across thousands of games. Statistical analysis of
these simulations reveals the strategic principles presented in this report.

\subsection{Organization}

Section~\ref{sec:rules} formalizes the game rules.
Section~\ref{sec:ai} details our AI algorithms.
Sections~\ref{sec:phase2} and~\ref{sec:phase1} analyze strategy for each phase.
Section~\ref{sec:results} presents simulation results.
Section~\ref{sec:conclusions} summarizes our findings.
